% !TeX spellcheck = en_US
\chapter{Introduction}\label{ch:introduction}%

This chapter motivates learning nut threading with a dexterous hand, defines the cyclic threading problem addressed in this thesis, states the research questions and contributions, and outlines the remaining chapters.

\section{Motivation}\label{sec:motivation}

Endowing robots with human-like dexterity is a long-standing goal of robotics, and one of its most convincing demonstrations is the ability to handle the parts and tools of a world built for human hands~\cite{BicchiKumar2000}. A multi-fingered hand promises to pick up objects of varying shape, hold and reorient them, and operate the very implements that people use every day. Learning-based methods have begun to realize this promise, and in-hand manipulation in particular has advanced rapidly in recent years~\cite{Weinberg2024,Chen2021,Handa2024}. Among the most demanding of these skills are those of mechanical assembly, where a firm grasp of a part must be coupled to a fine, contact-rich motion that joins it to another part.

Threading a nut onto a bolt is one of the most common operations in industrial assembly and, at the same time, one of the hardest to automate. The parts mate through sub-millimeter thread geometry, progress is governed by friction and contact forces that are difficult to sense, and success couples a stable grasp of the fastener to a controlled turning motion of the grasped part. Reinforcement learning (RL) in high-fidelity GPU simulation has recently made contact-rich assembly learnable: Factory~\cite{Narang2022} simulates nut-and-bolt engagement with thousands of simultaneous contacts in real time, and FORGE~\cite{Noseworthy2025} and IndustReal~\cite{Tang2023industreal} transfer such policies to physical robots. These systems, however, act with a parallel-jaw gripper and begin each episode with a fastener that is already held. The policy only learns to turn a pre-grasped part; grasp acquisition, grasp maintenance under turning, and the re-grasping that a full thread run-down requires are left outside the learning problem.

A dexterous hand changes the nature of the task. Humans thread nuts in several ways: a single finger suffices for free spinning, a pinch grasp is preferred for controlled threading, and a tripod or power grasp is applied under high resistance~\cite{Feix2016}. With a pinch grasp, one turn is bounded by the range of the wrist and the fingers, so a nut has to be turned, released, re-approached, and grasped again, cycle after cycle, while the nut advances along the bolt by one pitch per revolution. The task is therefore long-horizon and cyclic. Its difficulty does not lie in any single phase but in the transitions between phases: a grasp that is slightly off the intended faces, a release that disturbs the nut, or a return that lands the fingers at the wrong height propagates into the next cycle and compounds over ten or more cycles. Prior learning-based work on dexterous threading has addressed the spinning motion of an already engaged fastener~\cite{Hsieh2026}, and goal-conditioned screwing of furniture parts~\cite{Lum2026}, but not controlled, goal-directed, multi-cycle threading of everyday nuts with learned grasp and re-grasp.

\section{Problem Setup}\label{sec:problem-setup}

This thesis studies the task of threading a hexagonal nut onto a vertically mounted bolt with a seven-degree-of-freedom Franka Research~3 (FR3) arm equipped with an Inspire RH56DFX five-fingered hand. The nut starts engaged on the bolt by a fraction of a turn. One threading cycle consists of approaching the nut, establishing a pinch grasp between thumb and index finger on two opposing faces, rotating the nut clockwise by a target angle of $60\degree$ while it advances along the thread, holding the target, releasing the nut, and returning the hand to a re-grasp pose. A trial requires ten such cycles in sequence, and the deployed policy must execute them from an aligned RGB-D camera and the robot's proprioception alone, without privileged object state, scripted controllers, or human assistance.

The skill is learned in simulation and transferred to the physical robot. Simulating threads faithfully is computationally expensive, because engaged threads produce thousands of simultaneous contacts~\cite{Narang2022}; the reinforcement-learning stage of this work therefore uses a simplified thread model that couples a passive revolute joint to a position-controlled prismatic joint while retaining the finger--nut contact forces. The simplification creates a sim-to-real gap in contact and perception that the transfer method must close.

Two properties of the task shape the method. First, the task is long-horizon and cyclic: a policy must sustain consecutive cycles and is sensitive to errors at phase boundaries. Second, the task is only partially observable at deployment: the nut pose that the simulator exposes to a privileged policy must be inferred from camera images by the deployed policy. The first property favors generating the training signal phase by phase, where each phase is easy to specify; the second favors a teacher--student scheme in which a privileged teacher is distilled into a sensor-based student~\cite{Singh2025,Oertel2026}.

\section{Research Questions}\label{sec:research-questions}

The exposé of this thesis set out two research questions. The first asked whether a force-aware policy that threads a nut onto a bolt with a dexterous hand--arm system can be learned with proximal policy optimization (PPO) in a privileged simulation without constraining the thread interaction. The second asked whether such a privileged teacher can be distilled into a vision-based student that estimates the object pose with the accuracy required for sub-millimeter thread interaction. Both questions were answered, but not in the form they were posed. Repeated attempts to learn all phases of the cyclic task end to end with a single RL policy on faithfully simulated threads did not produce a threading teacher, because a reward that captures both long-horizon progress and precise transitions between phases could not be designed, and because the unconstrained thread contact admitted exploits such as pressing the nut through the thread. The thesis therefore refines the questions into the four that organize its experiments:

\begin{enumerate}
  \item Does a single student trained on complete multi-cycle demonstrations outperform repeated execution of a single-cycle policy (Q1)?
  \item Which visual representation preserves the geometry needed for repeated re-grasping under viewpoint and axial change, a point-cloud encoder or an image-space encoder (Q2)?
  \item Does a generative flow-matching action head distill the heterogeneous demonstrations more effectively than deterministic regression with recurrent or Transformer backbones (Q3)?
  \item Does augmenting simulated demonstrations with sparse human corrections collected during real-world replay close the sim-to-real gap (Q4)?
\end{enumerate}

\section{Contributions}\label{sec:contributions}

The main contributions of this thesis are the following.

\begin{enumerate}
  \item A formulation of long-horizon cyclic nut threading as an ordered sequence of phases with explicit cyclic and episodic tracking errors, and a hybrid teacher that combines a privileged RL policy for approach, grasp, rotation, and goal reaching with scripted release-and-return controllers. The hybrid teacher generates coherent multi-cycle demonstrations without a monolithic RL teacher.
  \item A human-corrected sim-to-real demonstration strategy in which simulated demonstrations are replayed on the physical robot and failure-prone segments between grasp and release are selectively overwritten by an operator, so that the demonstration distribution adapts to hardware contact and perception discrepancies.
  \item A unified visuomotor student distilled by conditional flow matching from point-cloud features of an aligned RGB-D view, proprioception, a global trajectory-progress scalar, and a process-phase encoding. The thesis derives the objective, documents the student as implemented, and places it against online DAgger variants, deterministic regression heads, alternative vision encoders, privileged observations, and two action contracts in a systematic ablation study.
  \item An experimental evaluation in simulation and on hardware. In ten-cycle trials the student completes every trial on the M24, M30, and M36 hexagonal nuts it was trained on and on unseen round and octagonal profiles, and the evaluation exposes rectangular profiles as a failure case of the learned re-grasp strategy.
\end{enumerate}

The method and its main results were submitted as a conference paper~\cite{NUTS2026}. This thesis expands the paper into a full report: it adds the background on imitation learning and flow matching, the design rationale for the action and observation contracts, the ablations implemented in the accompanying code base, the negative results that motivated the final design, and a discussion of the limitations.
% TODO(citation): add the NUTS paper itself to literature.bib once its venue and metadata are final.

\section{Outline}\label{sec:outline}

Chapter~\ref{ch:sota} reviews the theoretical background on reinforcement learning, imitation learning, generative action heads, visual representations, and operational-space control, and surveys related work on robotic assembly, policy chaining, sim-to-real transfer, and visuomotor distillation. Chapter~\ref{ch:method} presents the methodology: the problem formulation and metrics, the simulated and physical system, the privileged teacher, the hybrid teacher and demonstration pipeline, the conditional flow-matching student, and the evaluation protocol. Chapter~\ref{ch:sim} reports the experiments in simulation, including the ablations and the negative results that shaped the design. Chapter~\ref{ch:real} describes the hardware setup, the human-corrected replay, and the real-world results across nut sizes and geometries. Chapter~\ref{ch:discussion} interprets the results with respect to the research questions and discusses the limitations. Chapter~\ref{ch:conclusion} concludes and outlines future work.
